\section{Ethics, safety, and resources}

\subsection{Ethical considerations}
This work studies mechanistic aspects of research workflow design rather than deploying autonomous systems in human-facing domains. The evaluation uses constructed cases without human subjects or personally identifiable information. No web scraping, unsolicited messaging, or interaction with external platforms is performed by the evaluation system.

\subsection{Safety considerations}
ORION's design explicitly separates authority promotion (Paper IV) from formulation reconstruction (Paper I). Paper I mechanisms can expose possible method failures but cannot authorize $M_t$ to rewrite itself. The reframe gate is mechanically bounded: only formulation and search-universe coordinates may trigger a reframe, and evidence/execution-only negative controls verify that those gates do not fire inappropriately.

\subsection{Resource usage and environmental impact}
Model inference constitutes the primary computational cost. Token counts are logged for every system and case; approximate total usage for the local deterministic suite is recorded in the evidence manifest. Cloud provider regional selection for any external run follows institutional carbon-efficiency guidelines where applicable.

\subsection{Dual-use considerations}
The epistemic-state machinery described here is general-purpose and could be applied to non-scientific domains. The publication focuses on scientific research workflows and includes falsifier tests that verify negative-control behavior. No military or dual-use applications are developed or evaluated in this work.
